\section*{Chapter \ref{ch:iv:system-design} --- System Design}

\begin{solution}{iq:iv:system-design:1}
$200\,000 \times 23\,400 \times 64 \approx 2.995 \times 10^{11}$ bytes, about 300 gigabytes a day uncompressed, some 75
terabytes a year of 250 days. Compression reduces it by a factor that is worth measuring on the data
rather than assuming.
\iqlookfor{a clean chain with units and a year's extrapolation.}
\end{solution}

\begin{solution}{iq:iv:system-design:2}
$10^6 \times 64 \times 8 = 512$ megabits a second of payload; with 42 bytes of Ethernet, IPv4 and UDP headers ($14 + 20 + 8$) per
message sent alone, 848 megabits. A 1-gigabit link is marginal at best (and bursts inside the second are higher); 10 gigabits leaves room. Batching
several messages per packet cuts the header overhead.
\iqlookfor{payload versus wire bytes, and bursts shorter than a second.}
\end{solution}

\begin{solution}{iq:iv:system-design:3}
Instruments and order types; peak and average message rates, and the number of participants and sessions; latency targets (median,
99th and 99.9th percentile); matching rules (price-time or pro-rata, auctions, self-trade prevention); what market data is published
(depth, full order-by-order); durability and failover requirements (what may be lost, how fast to recover); regulatory records and
clock requirements.
\iqlookfor{requirements before design, including tails and failure.}
\end{solution}

\begin{solution}{iq:iv:system-design:4}
Add a limit order, cancel by identifier, modify (usually cancel and replace, losing priority unless only the size falls), match an
incoming order, and read the best price and the depth. With a sorted map of levels, a FIFO list per level and a hash index by
identifier: add at a new level $O(\log L)$ and at an existing one $O(1)$, cancel $O(1)$, best price $O(1)$, matching $O(1)$ per fill
(\cref{ch:iv:algorithms-and-data-structures}).
\iqlookfor{the operation set and the data-structure trio behind the complexities.}
\end{solution}

\begin{solution}{iq:iv:system-design:5}
Orders: $5 \times 10^6 \times 48 = 240$ megabytes. Levels: $10\,000 \times 2 \times 1\,000 \times 16 = 320$ megabytes. About 560
megabytes (534 MiB), which fits in memory easily but not in cache: the design question becomes locality (keep each instrument's hot
levels together, pre-allocate order records in pools) rather than capacity.
\iqlookfor{the arithmetic and the conclusion that locality, not size, matters.}
\end{solution}

\begin{solution}{iq:iv:system-design:6}
Sequence every inbound message and journal the sequenced input durably before acknowledging (or replicate it to a standby that
acknowledges); on restart, load the last snapshot and replay the journal from it. Gateways retransmit unacknowledged messages with
their client identifiers, and the engine discards duplicates it has already sequenced. It works only if the engine is deterministic:
its output is a function of the sequenced input alone, with no wall-clock reads, no iteration over unordered containers, and no
concurrency inside a partition (One Quant Book~13, chapter~24).
\iqlookfor{sequencing, journaling, replay, duplicate suppression and determinism.}
\end{solution}

\begin{solution}{iq:iv:system-design:7}
One publisher sequences updates per channel and sends them by multicast or, inside the firm, through a shared-memory or messaging
bus; each consumer has its own queue so that a slow one never blocks the publisher. A consumer that detects a gap requests a
retransmission or joins the snapshot channel and resumes from the snapshot's sequence number; a consumer too far behind is
disconnected and must resynchronise (\cref{fig:iv:system-design:md}). Monitor per-consumer lag.
\iqlookfor{sequencing, isolation of slow consumers, and a recovery path.}
\end{solution}

\begin{solution}{iq:iv:system-design:8}
An event loop over the merged, time-ordered recorded events; a simulated clock that the strategy reads through the same interface as
in production; latencies applied to data and to orders; a fill model that respects queue position and the recorded trades; fees and
rebates. Reproducible: data versions, code versions and random seeds recorded with each run. Believable: validated against live
fills of the same strategy (simulation--production parity, One Quant Book~15, chapter~12), with the fill model's assumptions stated.
\iqlookfor{event-driven structure, fill-model realism, and reproducibility.}
\end{solution}

\begin{solution}{iq:iv:system-design:9}
\Cref{lst:iv:system-design:match}: while the incoming order has quantity and the best opposite level crosses its limit, fill against the
level's orders in time order at the resting price, removing filled makers from the queue and the index and empty levels from the
map; any remainder rests at its limit behind earlier orders. Test it against a deliberately naive book (scan every resting order for
the best price and earliest time) on random order streams, comparing every fill and the best prices after each event, as the chapter's
C++ test does for 200 streams, and check invariants (the resting book never crossed).
\iqlookfor{correct price-time logic and a differential test against a brute force.}
\end{solution}

\begin{solution}{iq:iv:system-design:10}
Each gateway holds the state it needs to check an order locally in microseconds (its share of limits, its open orders and fills),
reserving against a limit before sending and releasing on cancel or fill; a central service holds the firm-wide positions and
limits, allocates limit shares to the gateways, and reconciles with the drop copies of the venues. If the central service is down,
gateways continue within their allocated shares and then stop: fail closed. A kill switch at each gateway and at the venues' cancel-on-
disconnect is independent of all of it (One Quant Book~11, chapter~27).
\iqlookfor{local fast state with central authority, reservation, and fail-closed behaviour.}
\end{solution}

\begin{solution}{iq:iv:system-design:11}
At half load a thread carries 500\,000 messages a second, so $3\,000\,000/500\,000 = 6$ partitions, about 1\,667 instruments each if
load were even. It is not: assign instruments by measured load, not by count, and keep hot instruments on their own partitions. An
order that spans partitions (a spread between two instruments) cannot be matched atomically by either: route such instruments to one
partition, or use a coordinator that reserves both legs, at a cost in latency.
\iqlookfor{the division, load balancing by activity, and the cross-partition problem.}
\end{solution}

\begin{solution}{iq:iv:system-design:12}
Run the standby as a replica consuming the same sequenced input (from the sequencer or the journal) in lock step, so that it holds
the same state at every sequence number. Acknowledge an order only when the sequenced message is on the standby as well (or in a
replicated journal), so no acknowledged order is lost. Detect the primary's death by heartbeats with a sub-second timeout; the standby
finishes processing its queued input, then takes over the gateways' sessions; clients reconnect and resend unacknowledged messages,
which the sequencer deduplicates. Fence the old primary so that it cannot resume after a partition.
\iqlookfor{a replicated state machine, the acknowledgement rule, fencing and deduplication.}
\end{solution}

\begin{solution}{iq:iv:system-design:13}
$100\,000 \times 50 \times 0.01 \times 1\,000 = 5 \times 10^7$ position updates a second. Do not revalue on every tick: index
positions by instrument so that a tick touches only its holders, revalue incrementally with sensitivities (delta and gamma) between
full revaluations, batch ticks into short intervals (the risk decision rarely needs every tick), and partition accounts across
workers; aggregate per account and per firm in trees.
\iqlookfor{the rate and the incremental, indexed, batched design that avoids it.}
\end{solution}
